Large language models trained with chain-of-thought prompting have shown
impressive multi-step reasoning, yet they remain fundamentally forward-only
reasoners: when a reasoning path fails, models cannot backtrack and
systematically explore alternatives.
We introduce \gore{} (\textbf{G}raph-\textbf{O}rchestrated \textbf{R}easoning
\textbf{E}nvironment), a minimal formal language that makes the search tree a
first-class syntactic object.
\gore{} has five primitives---\textsc{Step}, \textsc{Fork}, \textsc{Cut},
\letstmt{}, and \callstmt{}---whose operational semantics are structurally
isomorphic to Lean\,/\,Coq tactics (\texttt{exact}$\to$\textsc{Step},
\texttt{cases}$\to$\textsc{Fork}, \texttt{contradiction}$\to$\textsc{Cut},
\texttt{have}$\to$\letstmt{}, \texttt{apply}$\to$\callstmt{}).
This isomorphism enables a natural three-stage training pipeline: formal proof
pretraining instills an inductive bias for tree-structured search;
\gore{} curriculum fine-tuning transfers that bias to the explicit backtracking
syntax; and Dr.\,GRPO reinforcement learning with dense interpreter reward
aligns the model on verifiable reasoning chains.
Crucially, the \gore{} interpreter acts as a Process Reward Model, supplying
per-step verification signals without any human annotation.
We report an SFT baseline, a dense-vs.-sparse reward ablation (Dr.\,GRPO
vs.\ GRPO), and zero-shot transfer to natural-language reasoning benchmarks.
